% === Fig 2: repeatability by construction (gen_fig2_icc.py) ===
\begin{figure}[t]
\centering
\includegraphics[width=\textwidth]{figures/fig2_icc.pdf}
\caption{\textbf{Short-term repeatability of $A$ under each event construction, with one participant omitted, under recording eligibility and under detector settings.}
\textbf{a}, ICC(A,1) between R1 sessions 1 and 2 with 95\% participant-bootstrap CI (5,000 draws) for the six constructions, in the joint
cohort of participants eligible under all six ($n$ = 376, floor 5); the unit is the recording-level $A$ of each session. Open squares, unfiltered; filled circles, centred;
orange, mean-speed events; blue, admitted mean-speed events; green, per-eye events. $\ddagger$, post hoc comparator not in the frozen plan: mean-speed events after removing those with
either per-eye peak velocity above 1,000\textdegree/s (same cohort; .711 [.642, .769]).
\textbf{b}, Path-specific contrasts $\Delta$ICC (end $-$ start; u, unfiltered; c, centred) with 95\% CI from paired participant-bootstrap draws. Filled diamonds, all participants ($n$ = 376);
open circles, the same contrasts with the participant who contributed most to the session differences under unfiltered mean-speed events omitted ($n$ = 375). The contrasts share endpoints and are not
an additive partition of reliability. With windows, peaks and classes unchanged, admission raised ICC from .374 to .693,
$\Delta$ICC = .319 [.003, .501], or .054 [$-$.007, .129] with that participant omitted;
the ceiling comparator reproduced the increase, and admitted minus ceiling was $-$.018 [$-$.090, .039]. \textbf{c}, Mean-speed detector under different eligibility, floors 5 and 10:
each construction's own cohort, the common cohort eligible both unfiltered and centred, and a count-matched random-exclusion comparator ($\dagger$)
that keeps, in each participant, session and class, the number of events retained by centring (mean of 200 repetitions; its dashed bars without caps show the
2.5--97.5 percentile range of repetitions, not a CI). Dotted lines separate the own-cohort row from the common-cohort comparisons. Right-hand text: contrasts with 95\% CI in the common cohort ($n$ = 377 at floor 5;
$n$ = 342 at floor 10). \textbf{d}, Post hoc detector-setting sensitivity: all R1 recordings re-extracted with one detector setting changed at a time (velocity threshold, Savitzky--Golay window, maximum
left/right pairing gap in samples), each in its own joint cohort (right-hand $n$; floor 5; point estimates). Unfiltered mean-speed ICC ranged from .226 to .703,
admitted from .661 to .706 and ceiling from .697 to .736; re-extraction with the frozen settings reproduced the analysed events exactly.}
\label{fig:icc}
\end{figure}
